\chapter{Гиперзаряд из Yukawa-инвариантности и аномалий}

Первый вычислительный тест проверяет, следуют ли относительные заряды
наблюдаемого поколения из структурных требований.

Обозначим
\begin{equation}
\begin{gathered}
Y(Q_L)=q,\quad Y(u_R)=u,\quad Y(d_R)=d,\quad Y(L_L)=\ell,\\
Y(e_R)=e,\quad Y(\nu_R)=n,\quad Y(H)=h.
\end{gathered}
\end{equation}
Yukawa covariance даёт
\begin{equation}
u=q+h,\qquad d=q-h,\qquad e=\ell-h,\qquad n=\ell+h.
\end{equation}
Локальные anomaly coefficients равны
\begin{align}
\mathcal A_{331}&=2q-u-d,\\
\mathcal A_{221}&=3q+\ell,\\
\mathcal A_{\mathrm{grav}^2 1}&=6q-3u-3d+2\ell-e-n,\\
\mathcal A_{111}&=6q^3-3u^3-3d^3+2\ell^3-e^3-n^3.
\end{align}

После \(\ell=-3q\) все локальные аномалии при наличии \(\nu_R\)
сокращаются для любых \(q,h\). Остаётся непрерывное смешение с \(B-L\).
Условие стерильности \(n=0\) даёт \(h=3q\), поэтому при нормировке
\(h=1/2\)
\begin{equation}
\boxed{(q,u,d,\ell,e,n,h)=
\left(\frac16,\frac23,-\frac13,-\frac12,-1,0,\frac12\right).}
\end{equation}
Это target-отношения Стандартной модели.

Если \(\nu_R\) отсутствует, mixed gravitational и cubic anomalies сами
вынуждают \(h=3q\). Чистая цветовая аномалия сокращается как
\(2-1-1=0\), а число слабых дублетов \(3+1=4\) чётно, поэтому Witten
anomaly отсутствует.

\begin{theorem}[Условный вывод относительных гиперзарядов]
Yukawa covariance, anomaly cancellation и либо \(Y(\nu_R)=0\), либо
отсутствие \(\nu_R\), фиксируют стандартные относительные гиперзаряды.
Общая нормировка \(U(1)\)-генератора остаётся свободной.
\end{theorem}

Результат не выводит набор multiplets, число поколений, finite algebra или
gauge couplings. Машинная проверка сохранена в
\path{s2t_v4_hypercharge_anomaly_gate_results.json}.